\section*{अध्याय \ref{ch:b1:taylor} --- टेलर सूत्र और अनंतस्पर्शी प्रसार}

\begin{solution}{exo:b1:taylor:1}
\[
\eu^{2x} = 1 + 2x + 2x^2 + \frac{4x^3}{3} + o(x^3);
\qquad
\ln(1 - x) = -x - \frac{x^2}{2} - \frac{x^3}{3} - \frac{x^4}{4} +
o(x^4);
\]
\[
\sqrt{1+x} = 1 + \frac x2 - \frac{x^2}{8} + \frac{x^3}{16} + o(x^3);
\qquad
\frac{1}{1 + x^2} = 1 - x^2 + x^4 - x^6 + o(x^6),
\]
अंतिम वाला गुणोत्तर प्रसार में $-x^2$ प्रतिस्थापित करके।
\end{solution}

\begin{solution}{exo:b1:taylor:2}
$\eu^x - 1 - x = \frac{x^2}{2} + o(x^2)$: सीमा $\dfrac12$।

$\cos x = 1 - \frac{x^2}{2} + \frac{x^4}{24} + o(x^4)$ और $\sqrt{1 - x^2} =
1 - \frac{x^2}{2} - \frac{x^4}{8} + o(x^4)$: अंतर $\frac{x^4}{24} +
\frac{x^4}{8} = \frac{x^4}{6} + o(x^4)$: सीमा $\dfrac16$।

$\ln(1+x) - \sin x = \bigl(x - \frac{x^2}{2}\bigr) - x + o(x^2) =
-\frac{x^2}{2} + o(x^2)$: सीमा $-\dfrac12$।
\end{solution}

\begin{solution}{exo:b1:taylor:3}
$\frac{1}{1+t^2} = 1 - t^2 + t^4 + o(t^5)$; $0$ से $x$ तक समाकलन
(\cref{met:b1:taylor:compute}~(5)):
\[
\arctan x = x - \frac{x^3}{3} + \frac{x^5}{5} + o(x^5)\ \ (\text{सम}o(x^6)\text{, विषमता से}).
\]
$(1 - t^2)^{-1/2} = 1 + \frac{t^2}{2} + \frac38 t^4 + o(t^5)$ ($\alpha =
-\frac12$, $x = -t^2$ के साथ द्विपद प्रसार: $\binom{-1/2}{2} =
\frac{(-\frac12)(-\frac32)}{2} = \frac38$); समाकलन करने पर:
\[
\arcsin x = x + \frac{x^3}{6} + \frac{3x^5}{40} + o(x^5) .
\]
\end{solution}

\begin{solution}{exo:b1:taylor:4}
$a = 0$ पर $\exp$ के लिए टेलर--लाग्रांज (\cref{thm:b1:taylor:lagrange}), $x
= 1$: $\intcc{0}{1}$ पर $(n+1)$-वाँ \omterm{def:b1:derivative:def}{अवकलज} $\eu^t \leq \eu < 3$ है, अतः
\[
\Bigl|\eu - \sum_{k=0}^{n} \frac{1}{k!}\Bigr| \leq
\frac{3}{(n+1)!} .
\]
$6$ यथार्थ दशमलव के लिए $\frac{3}{(n+1)!} < 5\times 10^{-7}$ चाहिए, अर्थात्
$(n+1)! > 6\times 10^{6}$: चूँकि $10! = 3\,628\,800$ और $11! =
39\,916\,800$, $n + 1 = 11$, अर्थात् $n = 10$ पर्याप्त है।
\end{solution}

\begin{solution}{exo:b1:taylor:5}
$n\ln\bigl(1 + \frac1n\bigr) = n\Bigl(\frac1n - \frac{1}{2n^2} +
\frac{1}{3n^3} + o\bigl(\frac{1}{n^3}\bigr)\Bigr) = 1 - \frac{1}{2n} +
\frac{1}{3n^2} + o\bigl(\frac{1}{n^2}\bigr)$। चरघातांकी लेने पर, $u =
-\frac{1}{2n} + \frac{1}{3n^2}$ और $\eu^u = 1 + u + \frac{u^2}2 + o(u^2)$ के
साथ:
\[
\Bigl(1 + \frac1n\Bigr)^n = \eu\cdot \eu^{u}
= \eu\Bigl(1 - \frac{1}{2n} + \frac{1}{3n^2} + \frac{1}{8n^2}
+ o\Bigl(\frac{1}{n^2}\Bigr)\Bigr)
= \eu\Bigl(1 - \frac{1}{2n} + \frac{11}{24n^2} +
o\Bigl(\frac{1}{n^2}\Bigr)\Bigr).
\]
सीमा $\eu$; और त्रुटि $\sim \dfrac{\eu}{2n}$ है: धीमी (हर दस गुना $n$ पर एक
अंक)।
\end{solution}

\begin{solution}{exo:b1:taylor:6}
$f(x) = x^2 - x^4 = x^2(1 + o(1))$: पहला पद $x^2$, $p = 2$ सम, गुणांक $> 0$:
अतः $0$ पर स्थानीय न्यूनतम (समग्र नहीं: $f(2) = -12$)।

$g(x) = x^3 + x^5$: पहला पद $x^3$, विषम $p$: कोई चरम नहीं; $g$ अपनी
(क्षैतिज) स्पर्श रेखा काट देता है: $0$ पर नति-परिवर्तन।

$a = 1$ पर $h = \exp$: $h(x) = \eu + \eu(x-1) + \frac{\eu}{2}(x-1)^2 +
o((x-1)^2)$; स्पर्श रेखा से अंतर स्थानीय रूप से $\frac{\eu}{2}(x-1)^2 +
o(\cdot) > 0$ है। समग्र रूप से: उत्तलता
(\cref{thm:b1:derivative:convexchar}~(3)) से सभी $x$ के लिए $\eu^x - \eu x
\geq 0$: आलेख हर स्पर्श रेखा के ऊपर है, और समता केवल संपर्क बिंदु पर।
\end{solution}

\begin{solution}{exo:b1:taylor:7}
$x \to +\infty$ के लिए:
\[
f(x) = x\Bigl(1 + \frac1x\Bigr)^{1/3}
= x\Bigl(1 + \frac{1}{3x} - \frac{1}{9x^2} +
o\Bigl(\frac{1}{x^2}\Bigr)\Bigr)
= x + \frac13 - \frac{1}{9x} + o\Bigl(\frac1x\Bigr):
\]
अनंतस्पर्शी रेखा $y = x + \frac13$, और $+\infty$ के पास वक्र उसके नीचे। जब
$x \to -\infty$, तब वही संगणना वैध है (घनमूल सभी वास्तविक संख्याओं के लिए
परिभाषित है, और $\frac1x \to 0$): वही अनंतस्पर्शी रेखा $y = x + \frac13$, पर
अब $-\frac{1}{9x} > 0$: अर्थात् वक्र रेखा के \emph{ऊपर}।
\end{solution}

\begin{solution}{exo:b1:taylor:8}
$u_n = \sqrt n\bigl(\sqrt{1 + \tfrac1n} - 1\bigr) = \sqrt
n\bigl(\frac{1}{2n} + o(\frac1n)\bigr) \sim \dfrac{1}{2\sqrt n}$.

$v_n = \ln\bigl(1 + \frac1n\bigr) \sim \dfrac 1n$.

$w_n$: $h = \frac1n \to 0$, $\sin h - \tan h = \bigl(h - \frac{h^3}{6}\bigr)
- \bigl(h + \frac{h^3}{3}\bigr) + o(h^3) = -\frac{h^3}{2} + o(h^3)$ के साथ,
अतः $w_n \sim -\dfrac{1}{2n^3}$।
\end{solution}

\begin{solution}{exo:b1:taylor:9}
$x$ के इर्द-गिर्द कोटि $1$ पर, दोनों ओर, टेलर--लाग्रांज:
\[
f(x + h) = f(x) + h f'(x) + R_+,\quad
f(x - h) = f(x) - h f'(x) + R_-,
\qquad \abs{R_\pm} \leq \frac{h^2}{2} \sup \abs{f''} .
\]
घटाने पर: $f(x+h) - f(x-h) = 2h f'(x) + (R_+ - R_-)$, अतः
\[
\abs{f'(x)} \leq \frac{\abs{f(x+h) - f(x-h)}}{2h} +
\frac{h}{2}\sup_{\intcc{x-h}{x+h}}\abs{f''} .
\]
समग्र परिबंधों के साथ: प्रत्येक $h > 0$ के लिए $\abs{f'(x)} \leq
\frac{M_0}{h} + \frac{M_2 h}{2}$। दायाँ पक्ष $h = \sqrt{2M_0/M_2}$ पर
न्यूनतम होता है (\omterm{def:b1:derivative:def}{अवकलज} शून्य), जिसका मान $\sqrt{2M_0M_2}$ है --- अतः
$\abs{f'} \leq \sqrt{2M_0M_2}$। (यदि $M_2 = 0$, तो $h \to \infty$ लीजिए: $f'
= 0$, जो संगत है।)
\end{solution}

\begin{solution}{exo:b1:taylor:10}
$\intoo{0}{\pi}$ पर: $0 < \sin u < u$, अतः $(u_n)$ कठोरतः ह्रासमान और
धनात्मक है, इसलिए अभिसारी; उसकी सीमा $\intcc{0}{\pi}$ में $\sin$ का अचल
बिंदु है, और $\sin \ell = \ell$ $\ell = 0$ पर बाध्य कर देता है (क्योंकि $x >
0$ के लिए $\sin x < x$)।
\begin{enumerate}
  \item $u \to 0$ होने पर $\sin u = u - \frac{u^3}{6} + o(u^3)$ का प्रयोग
  करते हुए:
        \[
        v_{n+1} - v_n = \frac{1}{\sin^2 u_n} - \frac{1}{u_n^2}
        = \frac{1}{u_n^2}\Bigl(\Bigl(1 - \frac{u_n^2}{6} +
        o(u_n^2)\Bigr)^{\!-2} - 1\Bigr)
        = \frac{1}{u_n^2}\Bigl(\frac{u_n^2}{3} + o(u_n^2)\Bigr)
        \longrightarrow \frac13 .
        \]
  \item \cref{exo:b1:seq:10}~(3) (अंतरों के लिए चेज़ारो) से $\frac{v_n}{n}
  \to \frac13$, अर्थात् $v_n \sim \frac n3$, अर्थात् $u_n^2 \sim \frac 3n$:
  और चूँकि $u_n > 0$,
        \[
        u_n \sim \sqrt{\frac{3}{n}} .
        \]
\end{enumerate}
\end{solution}

\begin{solution}{exo:b1:taylor:11}
उभयनिष्ठ हर। पहली सीमा:
\[
\frac{1}{x^2} - \frac{1}{\sin^2 x}
= \frac{\sin^2 x - x^2}{x^2\sin^2 x},
\qquad
\sin^2 x = \Bigl(x - \frac{x^3}{6} + o(x^4)\Bigr)^{\!2}
= x^2 - \frac{x^4}{3} + o(x^5) ,
\]
अतः अंश $-\frac{x^4}{3} + o(x^4)$ है जबकि हर $\sim x^4$: इसलिए सीमा
$-\dfrac13$ है।

दूसरी: $\dfrac1x - \dfrac{1}{\ln(1+x)} = \dfrac{\ln(1+x) - x}{x\ln(1+x)}$;
अंश $-\frac{x^2}{2} + o(x^2)$ है और हर $x\bigl(x + o(x)\bigr) \sim x^2$: अतः
सीमा $-\dfrac12$ है।
\end{solution}

\begin{solution}{exo:b1:taylor:12}
$I_k = \intoo{k\pi - \frac\pi2}{k\pi + \frac\pi2}$ पर फलन $g(x) = \tan x -
x$ का \omterm{def:b1:derivative:def}{अवकलज} $\tan^2 x \geq 0$ है, जो केवल एक ही बिंदु $k\pi$ पर लुप्त होता
है: अतः $g$ $I_k$ पर कठोरतः वर्धमान है
(\cref{cor:b1:derivative:monotone}~(2)), और सिरों पर उसकी सीमाएँ $-\infty$
तथा $+\infty$ हैं: इसलिए ठीक एक शून्य $x_k$। $k \geq 1$ के लिए $g(k\pi) =
-k\pi < 0$, अतः $x_k \in \intoo{k\pi}{k\pi + \frac\pi2}$: $\varepsilon_k \in
\intoo{0}{\frac\pi2}$ के साथ $x_k = k\pi + \frac\pi2 - \varepsilon_k$ लिखिए।
तब
\[
x_k = \tan x_k = \tan\Bigl(\frac\pi2 - \varepsilon_k\Bigr)
= \frac{1}{\tan\varepsilon_k}
\quad\Longrightarrow\quad
\varepsilon_k = \arctan\frac{1}{x_k} ,
\]
जहाँ $\tan\varepsilon_k = \frac{1}{x_k}$ और $\varepsilon_k \in
\intoo{0}{\frac\pi2}$ प्रयुक्त हुए। चूँकि $x_k \geq k\pi \to \infty$:
$\varepsilon_k \to 0$, और
\[
\varepsilon_k = \arctan\frac{1}{x_k} = \frac{1}{x_k} +
O\Bigl(\frac{1}{x_k^3}\Bigr)
= \frac{1}{k\pi + O(1)} + O\Bigl(\frac{1}{k^3}\Bigr)
= \frac{1}{k\pi} + O\Bigl(\frac{1}{k^2}\Bigr) ,
\]
जिससे $x_k = k\pi + \frac\pi2 - \frac{1}{k\pi} + o\bigl(\frac1k\bigr)$।
\end{solution}

\begin{solution}{pb:b1:taylor:1}
\textbf{1.} $S_{2n+1} - S_{2n-1} = a_{2n} - a_{2n+1} \geq 0$ और $S_{2n+2} -
S_{2n} = a_{2n+2} - a_{2n+1} \leq 0$, जबकि $S_{2n} - S_{2n+1} = a_{2n+1} \to
0$: अतः अनुक्रम $(S_{2n+1})$, $(S_{2n})$ संलग्न हैं और उभयनिष्ठ $S$
(\cref{thm:b1:seq:adjacent}) की ओर अभिसरित होते हैं, जहाँ $S_{2n+1} \leq S
\leq S_{2n}$। सम $n$ के लिए: $S_{n+1} \leq S \leq S_n$ $-a_{n+1} \leq S -
S_n \leq 0$ देता है; विषम $n$ के लिए: $0 \leq S - S_n \leq a_{n+1}$। दोनों
स्थितियों में $\abs{S - S_n} \leq a_{n+1}$, और $S - S_n$ का चिह्न पहले छोड़े
गए पद $(-1)^{n+1}$ का चिह्न है। कठोर कमी हर दिखाई गई असमिका को कठोर बना देती
है, विशेष रूप से $0 < \abs{S - S_n} < a_{n+1}$।

\textbf{2.} $a_k = \frac{1}{k!}$ कठोरतः घटकर $0$ तक जाता है: अतः प्रश्न 1
लागू होता है। $-1$ और $0$ के बीच $\exp$ के लिए टेलर--लाग्रांज
(\cref{thm:b1:taylor:lagrange}): $\abs{\eu^{-1} - T_n} \leq
\frac{1}{(n+1)!}$ (वहाँ \omterm{def:b1:derivative:def}{अवकलज} $\eu^t$ $\leq 1$ है), अतः $T_n \to \eu^{-1}$,
और प्रश्न 1 की सीमा $S$ ठीक $\eu^{-1}$ \emph{है}, कठोर परिबंधों $0 <
\abs{\eu^{-1} - T_n} < \frac{1}{(n+1)!}$ के साथ।

\textbf{3.} सर्वसमिका का $\intcc{0}{1}$ पर समाकलन: बायाँ पक्ष $\arctan 1 =
\frac\pi4$ है (मूल प्रमेय), $k$-वाँ पद $\frac{(-1)^k}{2k+1}$ देता है, और
\[
\rho_n = (-1)^{n+1}\int_0^1 \frac{t^{2n+2}}{1+t^2}\dd t,
\qquad
\abs{\rho_n} \leq \int_0^1 t^{2n+2}\dd t = \frac{1}{2n+3} .
\]

\textbf{4.} $\pi$ पर त्रुटि $4\abs{\rho_n} \leq \frac{4}{2n+3}$ है:
$10^{-6}$ से नीचे लाने के लिए $2n + 3 > 4\cdot10^6$ चाहिए, अर्थात् लगभग बीस
लाख पद। इस बीच $4S_4 = 4\bigl(1 - \frac13 + \frac15 - \frac17 +
\frac19\bigr) = 4 \times 0.834921 = 3.339683$, जो $\pi$ से लगभग $0.2$ दूर
है: पाँच पद, और एक अंक भी नहीं।

\textbf{5.} $\intcc{0}{1}$ पर समाकलन: $R_n = \int_0^1 \frac{t^n}{1+t}\dd t$
के साथ $\ln 2 = \sum_{k=1}^{n} \frac{(-1)^{k-1}}{k} + (-1)^n R_n$; $\frac12
\leq \frac{1}{1+t} \leq 1$ से: $\frac{1}{2(n+1)} \leq R_n \leq
\frac{1}{n+1}$। त्रुटि $\frac1n$ के दो गुणजों के बीच फँसी है: छह दशमलव की
क़ीमत लगभग दस लाख पद।

\textbf{6.} $0$ से $x$ तक समाकलन: $\frac12\ln\frac{1 + x}{1 - x} =
\sum_{k=0}^{n} \frac{x^{2k+1}}{2k+1} + \int_0^x \frac{t^{2n+2}}{1-t^2}\dd
t$। $x = \frac13$ पर: $\frac{1 + 1/3}{1 - 1/3} = 2$, और $\intcc{0}{\frac13}$
पर $\frac{1}{1 - t^2} \leq \frac98$:
\[
\ln 2 = 2\sum_{k=0}^{n} \frac{(1/3)^{2k+1}}{2k+1} +
\tilde\rho_n,
\qquad
0 < \tilde\rho_n \leq \frac{9}{4}\cdot
\frac{(1/3)^{2n+3}}{2n+3} :
\]
हर अतिरिक्त पद त्रुटि को लगभग $9$ से भाग देता है।

\textbf{7.} आगमन: $n = 1$ के लिए $1 - \frac12 = \frac12 = H_2 - H_1$। पद:
\[
H_{2n+2} - H_{n+1} = (H_{2n} - H_n) + \frac{1}{2n+1} +
\frac{1}{2n+2} - \frac{1}{n+1}
= (H_{2n} - H_n) + \frac{1}{2n+1} - \frac{1}{2n+2},
\]
जो ठीक एकांतरित योग की वृद्धि है। और $H_{2n} - H_n =
\sum_{k=1}^{n}\frac{1}{n+k}$ \cref{ex:b1:integration:riemannexample} का
\omterm{thm:b1:integration:riemann}{रीमान योग} है, जो $\ln 2$ की ओर अभिसरित होता है: अर्थात् तरीक़ा 1 के सम आंशिक
योग तरीक़ा 3 के \omterm{thm:b1:integration:riemann}{रीमान योग} \emph{ही} हैं।

\textbf{8.} $\sum_{k=1}^{6}\frac{(-1)^{k-1}}{k} = 0.61667$, त्रुटि $0.0765$;
और $n = 5$ पर तरीक़ा 2 $\leq \frac94\cdot\frac{(1/3)^{13}}{13} =
1.1\cdot10^{-7}$ त्रुटि के साथ $0.6931471$ देता है। कारण: तरीक़ा 1 लघुगणक
श्रेणी का मान सीमांत बिंदु $x = 1$ पर लेता है, जहाँ पद $\frac1k$ की तरह क्षय
होते हैं; तरीक़ा 2 उसका मान $x = \frac13$ पर लेता है, जो भीतर गहरा है और
जहाँ हर पद एक नया गुणक $\frac19$ ढोता है।

\textbf{9.} दस दशमलव: $\tilde\rho_n \leq 5\cdot10^{-11}$ चाहिए। $n = 10$ पर:
$\frac94 \cdot \frac{(1/3)^{23}}{23} =
\frac94\cdot\frac{1.06\cdot10^{-11}}{23} \approx 1.0\cdot10^{-12} <
5\cdot10^{-11}$: अतः ग्यारह पद पर्याप्त हैं।

\textbf{10.} $(5+\iu)^2 = 24 + 10\iu$, फिर $(5+\iu)^4 = (24 + 10\iu)^2 = 476
+ 480\iu$; और $2(1+\iu)(239+\iu) = 2(238 + 240\iu) = 476 + 480\iu$: बराबर।
कोणांक: $\arg(5 + \iu) = \arctan\frac15$, अतः बाएँ पक्ष का कोणांक
$4\arctan\frac15 \approx 0.79 \in \intoo{0}{\pi}$ है; और दाएँ पक्ष का कोणांक
$\frac\pi4 + \arctan\frac{1}{239} \in \intoo{0}{\pi}$ है। एक ही $< 2\pi$
लंबाई के \omterm{prop:b1:reals:intervals}{अंतराल} में कोणांक वाली दो बराबर सम्मिश्र संख्याएँ:
\[
4\arctan\frac15 = \frac\pi4 + \arctan\frac{1}{239} ,
\]
जो मैकिन का सूत्र है।

\textbf{11.} $\frac{1}{1+t^2} = \sum_{k=0}^n (-1)^k t^{2k} +
\frac{(-1)^{n+1}t^{2n+2}}{1+t^2}$ का $0$ से $x$ तक समाकलन कीजिए:
\[
\arctan x = \sum_{k=0}^{n}\frac{(-1)^k x^{2k+1}}{2k+1} +
r_n(x),
\qquad
\abs{r_n(x)} \leq \int_0^x t^{2n+2}\dd t =
\frac{x^{2n+3}}{2n+3} .
\]

\textbf{12.} त्रुटियाँ: $16\,\frac{(1/5)^{11}}{11} = 3.0\cdot 10^{-8}$ और
$4\,\frac{(1/239)^5}{5} \approx 1\cdot10^{-12}$: कुल $< 5\cdot10^{-8}$।
दिखाया गया योग $3.14159268\dots$ निकलता है, अतः $\pi = 3.1415926\dots$
$5\cdot10^{-8}$ तक प्रमाणित: अर्थात् छह पदों से सात दशमलव (उदारता से गिनें
तो $\frac15$ पर पाँच और $\frac1{239}$ पर दो)।

\textbf{13.} लाइब्निज़: त्रुटि $\sim \frac1n$, अतः हर नया अंक काम को दस गुना
कर देता है। डाल्ज़ेल (\cref{pb:b1:integration:1}, प्रश्न 22): त्रुटि
$4^{1-5m}$, अर्थात् लगभग तीन अंक प्रति पद, पर हर पद में भारी \omterm{def:b1:poly:def}{बहुपद}। मैकिन:
प्रति पद त्रुटि-अनुपात $\frac{1}{25}$, अर्थात् लगभग $1.4$ अंक प्रति पद, और
हर पद में एक ही भाग। उपदेश: गुणोत्तर प्रकार के प्रसार का शेषफल $x^{2n}$ की
तरह मापित होता है, अतः $x$ को छोटा करना प्रति पद \emph{नियत} लागत पर अंक
ख़रीद लेता है --- और मैकिन की सम्मिश्र सर्वसमिका ठीक $x$ सिकोड़ने की मशीन
है।

\textbf{14.} $a_k = \frac{1}{(2k)!}$ कठोरतः घटकर $0$ तक जाता है; प्रश्न 1 और
\cref{prop:b1:taylor:standard} (प्रश्न 2 की भाँति लाग्रांज परिबंध) से
$\sum_{k \leq n}\frac{(-1)^k}{(2k)!} \to \cos 1$, और वह भी \emph{कठोर} घेराव
$0 < \bigl|\cos 1 - S'_n\bigr| < \frac{1}{(2n+2)!}$ के साथ। मान लीजिए $\cos
1 = \frac pq$ और $2n \geq q$ लीजिए: तब $(2n)!\,S'_n = \sum_{k\leq n} (-1)^k
\frac{(2n)!}{(2k)!} \in \Z$ और $(2n)!\,\frac pq \in \Z$, जबकि
\[
0 < \Bigl|(2n)!\cos 1 - (2n)!S'_n\Bigr| <
\frac{(2n)!}{(2n+2)!} = \frac{1}{(2n+1)(2n+2)} < 1 :
\]
अर्थात् $< 1$ निरपेक्ष मान वाला अशून्य पूर्णांक। विरोधाभास: $\cos 1 \notin
\Q$।

\textbf{15.} ठीक वैसे ही $a_k = \frac{1}{(2k+1)!}$ के साथ, $2n + 1 \geq q$
वाले $(2n+1)!$ से गुणा करने पर: $\sin 1 \notin \Q$। फिर भी $\cos^2 1 +
\sin^2 1 = 1 \in \Q$: अपरिमेय संख्याओं के गुणनफल और योग परिमेय हो सकते हैं
--- अपरिमेयता किसी भी बीजीय संक्रिया से मुफ़्त में पार नहीं होती।

\textbf{16.} पूँछ-परिबंध: $m > n$ के लिए,
\[
\sum_{k=n+1}^{m} \frac{1}{(2k)!}
\leq \frac{1}{(2n+2)!}\Bigl(1 + \frac12 + \frac14 +
\dots\Bigr) \leq \frac{2}{(2n+2)!} ,
\]
क्योंकि हर क्रमिक अनुपात $\frac{1}{(2k+1)(2k+2)} \leq \frac12$ है; और पूँछ
धनात्मक है (उसका पहला पद धनात्मक है)। अतः $0 < \cosh 1 - \sum_{k\leq
n}\frac{1}{(2k)!} < \frac{2}{(2n+2)!}$, और $2n \geq q$ वाले $(2n)!$ से गुणा
करने पर फिर एक अशून्य पूर्णांक $\intoo{0}{1}$ में फँस जाता है: $\cosh 1
\notin \Q$।

\textbf{17.} $\cos\frac1m = \sum_k \frac{(-1)^k}{m^{2k}(2k)!}$: पद कठोरतः
घटकर शून्य तक जाते हैं, और $k \leq n$ के लिए
$m^{2n}(2n)!\cdot\frac{1}{m^{2k}(2k)!} = m^{2(n-k)}\frac{(2n)!}{(2k)!} \in
\Z$। यदि $\cos\frac1m = \frac pq$, तो कठोर घेराव को $q\,m^{2n}(2n)!$ से गुणा
कीजिए: बड़े $n$ के लिए त्रुटि $\frac{q}{m^2(2n+1)(2n+2)} < 1$ से परिबद्ध है:
विरोधाभास। $a \geq 2$ वाले $\frac ab$ के लिए: हर हटाने से पूँछ $b^{2n}(2n)!$
से गुणित हो जाती है, पर पहला छोड़ा गया पद $\frac{a^{2n+2}}{b^{2n+2}(2n+2)!}$
है, और गुणनफल $\frac{a^{2n+2}}{b^2(2n+1)(2n+2)}$ \emph{फट पड़ता है}: अंश
$a^{2n+2}$ अब नहीं कटता, और जाल जाम हो जाता है। (परिणाम फिर भी सत्य है ---
पर नाइवन-शैली की मशीनरी से, इससे नहीं।)

\textbf{18.} $t = x$ पर $g$ का हर पद $f(x) - f(x) = 0$ को छोड़कर लुप्त हो
जाता है: $g(x) = 0$; और $A$ ऐसा चुना गया है कि $g(a) = 0$। अवकलन करने पर योग
दूरबीनी हो जाता है:
\[
g'(t) = -\frac{f^{(n+1)}(t)}{n!}(x - t)^n +
A\,\frac{(x-t)^n}{n!}
= \frac{(x-t)^n}{n!}\bigl(A - f^{(n+1)}(t)\bigr) .
\]
$a$ से $x$ तक के खंड पर रोल कठोरतः बीच में ऐसा $c$ देता है कि $g'(c) = 0$;
और चूँकि $(x - c)^n \neq 0$: $A = f^{(n+1)}(c)$। $g(a) = 0$ को खोलने पर
शेषफल $\frac{f^{(n+1)}(c)}{(n+1)!}(x-a)^{n+1}$ वाली टेलर समता मिल जाती है।

\textbf{19.} $x > 0$ के लिए शेषफल $\frac{\eu^{c}} {(n+1)!}x^{n+1} > 0$ है:
अर्थात् चरघातांकी हर कोटि पर अपने प्रत्येक टेलर \omterm{def:b1:poly:def}{बहुपद} से कठोरतः बड़ा है। $x
< 0$ के लिए शेषफल का चिह्न $x^{n+1}$ का चिह्न है: $\eu^x$ विषम $n$ के लिए
\omterm{def:b1:poly:def}{बहुपद} के \emph{ऊपर} है और सम $n$ के लिए \emph{नीचे} --- अर्थात् बारी-बारी से
दोनों ओर, जैसा $\eu^x$ के सामने $1 + x$ और $1 + x + \frac{x^2}{2}$ के आलेख
पहले ही दिखा देते हैं।

\textbf{20.} सच्ची त्रुटि: $\sin 0.5 - 0.4791667 = 2.59\cdot10^{-4}$, परिबंध
$\frac{0.5^5}{120} = 2.60\cdot10^{-4}$ के सामने: अर्थात् लगभग प्राप्त (अगला
पद पूँछ पर हावी है)। यंग: सीमाओं और स्थानीय विश्लेषण के लिए, जहाँ किसी अचर
की आवश्यकता नहीं। लाग्रांज: प्रमाणित दशमलवों के लिए। एकांतरित: जहाँ लागू हो,
वही परिबंध \emph{और साथ में} त्रुटि की दिशा --- तीनों में सर्वोत्तम, पर सबसे
दुर्लभ।

\textbf{21.} $0$ पर \omterm{def:b1:continuity:continuous}{सांतत्य}: $u = \frac{1}{x^2} \to +\infty$ के साथ $f(x) =
\eu^{-u} \to 0 = f(0)$। $0$ पर \omterm{def:b1:derivative:def}{अवकलज}: $\bigl|\frac{f(h)}{h}\bigr| = \sqrt
u\,\eu^{-u} \to 0$ (\cref{prop:b1:functions:powerrules}): $f'(0) = 0$। $x
\neq 0$ के लिए $f'(x) = \frac{2}{x^3}\eu^{-1/x^2}$: अर्थात् $P_1(X) = 2X^3$
के साथ $P_1\bigl(\frac1x\bigr)\eu^{-1/x^2}$ रूप; और आगमन से
$P_k(\frac1x)\eu^{-1/x^2}$ का अवकलन $P_{k+1}(X) = 2X^3 P_k(X) - X^2 P_k'(X)$
देता है, जो \omterm{def:b1:poly:def}{बहुपद} है। फिर
\[
\frac{f^{(k)}(h) - 0}{h} = v\,P_k(v)\,\eu^{-v^2}
\Big|_{v = 1/h} \longrightarrow 0
\]
($\pm\infty$ पर \omterm{prop:b1:functions:powerrules}{वृद्धि तुलना}, $\eu^{-v^2}$ के सामने \omterm{def:b1:poly:def}{बहुपद}): अतः आगमन से सभी
$k$ के लिए $f^{(k)}(0) = 0$। $0$ पर $f$ के सभी टेलर \omterm{def:b1:poly:def}{बहुपद} लुप्त हो जाते हैं,
फिर भी $0$ से बाहर $f > 0$: अर्थात् टेलर--यंग हर कोटि पर यथार्थ है और अंकुर
से आगे अंधा। प्रसार केवल स्थानीय सूचना है।

\textbf{22.} प्रश्न 2 से कठोरतः $0 < \abs{\eu^{-1} - T_n} <
\frac{1}{(n+1)!}$। यदि $\eu^{-1} = \frac pq$, तो $n \geq q$ लीजिए और $n!$ से
गुणा कीजिए: $n!\,T_n \in \Z$ और $n!\frac pq \in \Z$, अतः किसी अशून्य
पूर्णांक का निरपेक्ष मान $< \frac{n!}{(n+1)!} = \frac{1}{n+1} < 1$ हो जाता
है: विरोधाभास। अतः $\eu^{-1} \notin \Q$, और $\eu = \frac{1}{\eu^{-1}}$ भी
अपरिमेय है। तीनों उपपत्तियाँ: \cref{exo:b1:seq:9} में $q!\,\eu$ को दबाने
वाले \omterm{thm:b1:seq:adjacent}{संलग्न अनुक्रम}; \cref{pb:b1:integration:1} में \omterm{thm:b1:integration:def}{समाकल} पुनरावृत्ति $A_n =
\eu - nA_{n-1}$; और एकांतरित घेराव (यहाँ)। एक जाल, लघुता के तीन प्रमाणपत्र।

\textbf{23.} आंशिक योगों को जोड़ों में समूहबद्ध कीजिए: $E_n = \sum_{j=1}^{n}
\frac{1}{4j^2}$ वाला $\sum_{k=1}^{2n} (-1)^k b_k = E_n - O_n$, जो परिबद्ध है
(\cref{ex:b1:seq:basel}, $E_n \leq \frac12$ के दूरबीनी परिबंध से), और $O_n =
\sum_{j=1}^{n}\frac{1}{2j-1} \geq \frac12 H_n \to +\infty$
(\cref{exo:b1:seq:5}): अतः आंशिक योग $-\infty$ की ओर जाते हैं। प्रश्न 1 में
एकदिष्टता ठीक वहाँ प्रयुक्त हुई थी जहाँ $S_{2n+1} - S_{2n-1} = a_{2n} -
a_{2n+1}$ को चिह्न चाहिए था: कमी के बिना सम और विषम उपानुक्रमों का एकदिष्ट
होना आवश्यक नहीं, और संलग्नता ढह जाती है।

\textbf{24.} $(2+\iu)(3+\iu) = 5 + 5\iu = 5(1+\iu)$; कोणांक लेने पर (सभी
$\intoo{0}{\frac\pi2}$ में): $\arctan\frac12 + \arctan\frac13 = \frac\pi4$।
छह दशमलव के लिए श्रेणी की लागत: त्रुटि $\leq
4\bigl(\frac{(1/2)^{2n+3}}{2n+3} + \frac{(1/3)^{2n+3}}{2n+3}\bigr)$, जो $n =
10$ पर $\approx 2\cdot10^{-8} < 5\cdot10^{-7}$ है: अर्थात् ग्यारह पद। क्रम:
लाइब्निज़ से चरघातांकी अंतर से बेहतर, और मैकिन से पीछे (जिसका प्रभावी बिंदु
$\frac15$ $\frac12$ से छोटा है): मैकिन के $1.4$ के सामने लगभग $0.6$ अंक
प्रति पद।

\textbf{25.} (क) ह्रासमान $a_k \to 0$ के लिए एकांतरित आंशिक योग $\abs{S -
S_n} \leq a_{n+1}$ के साथ अभिसरित होते हैं, और त्रुटि पहले छोड़े गए पद का
चिह्न ढोती है। (ख) स्पष्ट शेषफल वाली परिमित सर्वसमिका का \emph{किसी चुने हुए
बिंदु पर मान निकाला और परिबद्ध किया} जा सकता है, जबकि सीमा-कथन केवल अंततः
निकटता का वचन देता है --- और प्रमाणन के लिए पहला ही चाहिए। (ग) निकाले गए
परिणाम: मैकिन से $5\cdot10^{-8}$ तक $\pi$, $\frac13$-श्रेणी से दस दशमलव तक
$\ln 2$, $\cos 1$, $\sin 1$, $\cosh 1$, $\cos\frac1m$ और $\eu^{-1}$ की
अपरिमेयता, टेलर--लाग्रांज समता, और चपटे फलन की चेतावनी। (घ)
\cref{ch:b1:series} प्रश्न 1 को एकांतरित श्रेणी कसौटी में उन्नत करता है,
निरपेक्ष को सप्रतिबंध अभिसरण से अलग करता है, और उसकी सप्ताहांत समस्या वह
पुनर्व्यवस्था-नाटक मंचित करती है जिसका मुख्य साक्षी तरीक़ा 1 की एकांतरित
हरात्मक श्रेणी है।
\end{solution}
